% Part: first-order-logic
% Chapter: models-theories
% Section: introduction

\documentclass[../../../include/open-logic-section]{subfiles}

\begin{document}

\olfileid[hi]{fol}{mat}{int}
\olsection{परिचय}

\begin{explain}
अभिगृहीतीय विधि का विकास विज्ञान के इतिहास की एक महत्त्वपूर्ण उपलब्धि है,
और गणित के इतिहास में उसका विशेष महत्त्व है। किसी क्षेत्र के अभिगृहीतीय
विकास में अनेक प्रश्नों को स्पष्ट करना पड़ता है: यह क्षेत्र किस विषय का
अध्ययन करता है? इसकी सबसे मूलभूत अवधारणाएँ कौन-सी हैं? वे आपस में किस प्रकार
संबंधित हैं? क्या क्षेत्र की सभी अवधारणाओं को इन मूलभूत अवधारणाओं के पदों में
परिभाषित किया जा सकता है? इन अवधारणाओं को किन नियमों का पालन करना पड़ता है,
और अनिवार्यतः करना चाहिए?

अभिगृहीतीय विधि और तर्कशास्त्र मानो एक-दूसरे के लिए ही बने हैं। औपचारिक
तर्कशास्त्र अभिगृहीतीय सिद्धांतों को सूत्रबद्ध करने, सिद्धांत के अभिगृहीतों से
ठीक-ठीक निर्दिष्ट रीति से प्रमेय सिद्ध करने, और अभिगृहीतों को संतुष्ट करने वाले
सभी तंत्रों के गुणों का व्यवस्थित अध्ययन करने के साधन देता है।
\end{explain}

\begin{defn}
!!{sentence}s का कोई समुच्चय~$\Gamma$ \emph{संवृत} है, तभी और केवल तभी जब
$\Gamma \Entails !A$ होने पर $!A \in \Gamma$। !!{sentence}s के समुच्चय
$\Gamma$ का \emph{संवरण} $\Setabs{!A}{\Gamma \Entails !A}$ है।

यदि~$\Gamma$, वाक्यों के समुच्चय~$\Delta$ का संवरण है, तो हम कहते हैं
कि~$\Gamma$,~$\Delta$ \emph{द्वारा अभिगृहीतीकृत} है।
\end{defn}

\begin{explain}
किसी अभिगृहीतीय सिद्धांत को हम उन वाक्यों का समुच्चय मान सकते हैं जो
उसके अभिगृहीत-समुच्चय~$\Delta$ द्वारा अभिगृहीतीकृत है। दूसरे शब्दों में, मान लें
कि हमारे पास उस अभिगृहीतीय रूप से विकसित विज्ञान की मूल अवधारणाओं के लिए
अतार्किक चिह्नों वाली प्रथम-क्रम भाषा और उस विज्ञान के मूलभूत नियमों को व्यक्त
करने वाला !!{sentence}s का समुच्चय है। तब सिद्धांत को इस भाषा के उन सभी
!!{sentence}s द्वारा निरूपित मान सकते हैं जो अभिगृहीतों से अर्थगत रूप से
निष्पन्न होते हैं। इसके उदाहरण किसी एक मूल अवधारणा और सरल अभिगृहीतों वाले
आंशिक क्रमों के सिद्धांत से लेकर न्यूटनियन यांत्रिकी जैसे जटिल सिद्धांतों तक
फैले हैं।

अभिगृहीतीय विधि के इस औपचारिक रूप को इतना महत्त्वपूर्ण बनाने वाले मुख्य
तार्किक तथ्य निम्नलिखित हैं। मान लें कि~$\Gamma$ किसी सिद्धांत का
अभिगृहीत-तंत्र, अर्थात् वाक्यों का समुच्चय है।
\begin{enumerate}
\item हम ठीक-ठीक बता सकते हैं कि कोई अभिगृहीत-तंत्र संरचनाओं के अभीष्ट वर्ग को
  कब निरूपित करता है; दूसरे शब्दों में, उसके प्रतिरूप कौन-सी !!{structure}s
  हैं। अर्थात्, यदि कोई निश्चित वर्ग अभीष्ट है और उसकी सदस्य !!{structure}s
  हैं, तो अभिगृहीत-तंत्र~$\Gamma$ उस वर्ग को तभी और केवल तभी ठीक निरूपित करेगा,
  जब उस वर्ग में आने वाली !!{structure}s ठीक वे~$\Struct M$ हों जिनके लिए
  $\Sat{M}{\Gamma}$ हो।
\item इस दृष्टि से हम असफल हो सकते हैं, क्योंकि ऐसे~$\Struct M$ हो सकते हैं
  जिनके लिए $\Sat{M}{\Gamma}$ हो, किंतु~$\Struct M$ अभीष्ट न हो; हमारे लिए
  अभीष्ट !!{structure}s दूसरी हों। तब हम ऐसे अभिगृहीत जोड़ सकते हैं जो~$\Struct M$
  में सत्य नहीं हैं।
\item यदि हम कम-से-कम इतना सफल हों कि जो भी !!{structure}s अभीष्ट हों, वे सभी
  $\Gamma$ को सत्य बनाती हों, तो वे सभी !!{structure}s वाक्य~$!A$ को भी सत्य
  बनाती हैं, जब भी $\Gamma \Entails !A$। अतः यह दिखाने के लिए कि सभी अभीष्ट
  !!{structure}s किसी वाक्य को सत्य बनाती हैं, हम केवल यह दिखा सकते हैं कि वह
  अभिगृहीतों से निष्पन्न होता है; इसके लिए तार्किक साधनों (जैसे
  !!{derivation} विधियों) का उपयोग किया जा सकता है।
\item कभी-कभी हमारे मन में पहले से कोई अभीष्ट !!{structure} नहीं होती और हम
  स्वयं अभिगृहीतों से आरंभ करते हैं: हम कुछ मूल अवधारणाएँ लेते हैं, जिनसे हम
  कुछ निश्चित नियमों का पालन कराना चाहते हैं, और उन नियमों को अभिगृहीत-तंत्र
  में संहिताबद्ध करते हैं। हम तुरंत यह जाँचना चाहेंगे कि अभिगृहीत परस्पर
  विरोधी नहीं हैं। यदि वे विरोधी हैं, तो इन नियमों को मानने वाली कोई अवधारणा
  हो ही नहीं सकती और हमने एक असंगत सिद्धांत बनाने का प्रयास किया है।
  $\Gamma$ का कोई प्रतिरूप खोजकर हम जाँच सकते हैं कि ऐसा नहीं होता। यदि हमारे
  सिद्धांत के प्रतिरूप हैं, तो हम तार्किक विधियों से उनका अध्ययन और निर्माण
  दोनों कर सकते हैं।
\item अभिगृहीतों की स्वतंत्रता भी एक महत्त्वपूर्ण प्रश्न है। हो सकता है कि
  कोई अभिगृहीत वास्तव में अन्य अभिगृहीतों का परिणाम हो और इसलिए अनावश्यक हो।
  अभिगृहीत~$!A$ को, जो~$\Gamma$ में है, अनावश्यक सिद्ध करने के लिए हम $\Gamma \setminus
  \{!A\} \Entails !A$ सिद्ध कर सकते हैं। यह दिखाकर कि $(\Gamma \setminus \{!A\}) \cup \{\lnot
  !A\}$ संतुष्टनीय है, हम यह भी सिद्ध कर
  सकते हैं कि अभिगृहीत अनावश्यक नहीं है। उदाहरणतः इसी प्रकार दिखाया गया कि
  समांतर अभिगृहीत, ज्यामिति के अन्य अभिगृहीतों से स्वतंत्र है।
\item एक और महत्त्वपूर्ण प्रश्न सिद्धांत की अवधारणाओं की परिभाषणीयता का है:
  भाषा का चुनाव निर्धारित करता है कि सिद्धांत के प्रतिरूप किन वस्तुओं और
  संबंधों से बने हैं। किंतु सिद्धांत के प्रत्येक पक्ष को उसके प्रतिरूपों में
  अलग से निरूपित करना आवश्यक नहीं। उदाहरणतः प्रत्येक क्रम~$\le$ एक तदनुरूप
  प्रबल क्रम~$<$ निर्धारित करता है---इनमें से एक दिया हो तो दूसरे को परिभाषित
  कर सकते हैं। अतः ऐसे क्रम वाले किसी सिद्धांत के प्रतिरूप में तदनुरूप प्रबल
  क्रम का \emph{अलग से} होना आवश्यक नहीं। सामान्यतः एक संबंध को अन्य संबंधों
  के पदों में कब परिभाषित किया जा सकता है? और किसी संबंध को अन्य संबंधों के
  पदों में परिभाषित करना कब असंभव है (जिससे उसे भाषा की मूल अवधारणाओं में
  जोड़ना पड़े)?
\end{enumerate}
\end{explain}


\end{document}
